\section[The A-hypergeometric system]{The ${\boldsymbol A}$-hypergeometric system}\label{chap-ahyp}

Fix a positive integer $N$ and let $\bm{\gamma}\in\mathbb{R}^N$ be a row vector. Let $L\subset\mathbb{Z}^N$ be a lattice of
rank $d$ which satisfies the following conditions:
\begin{enumerate}\itemsep=0pt
\item $L$ is contained in the hyperplane $\big((l_1,l_2,\dots,l_N) \in \mathbb{Z}^N, \, \sum_{i=1}^Nl_i=0\big)$.
\item $L$ is saturated, i.e., $(L\otimes\mathbb{R})\cap\mathbb{Z}^N=L$.
\end{enumerate}
Now define
\begin{gather*}
\Phi_{\bm{\gamma}}^L:=\sum_{\bm{l}\in L}\prod_{j=1}^N\frac{z_j^{\gamma_j+l_j}}
{\Gamma(\gamma_j+l_j+1)}.
\end{gather*}
For the moment this is a formal series expansion. Throughout this paper we denote (row)-vectors with a bold face and its components have a normal font weight (e.g., $\bm{\gamma} = (\gamma_1,\dots,\gamma_N )$).
Notice that $\Phi_{\bm{\gamma}}^L
=\Phi_{\bm{\gamma}+\bm{l}}^L$ for any $\bm{l}\in L$.
Let $r = N-d$ and let $A$ be an $(r \times N)$-matrix with integer entries such that $L$ is
the integer kernel of $A$.
Let us define
$\bm{\alpha}=A\bm{\gamma}^\intercal$. Notice that $A(\bm{\gamma}+\bm{l})^\intercal
=\bm{\alpha}$ for any $\bm{l}\in L\otimes\mathbb{R}$.
We call this the parameter vector of the $A$-hypergeometric system we will define.
Because $L$ is contained in the hyperplane $\sum_{i=1}^Nl_i=0$, there is a linear form $h\colon \mathbb{R}^{r} \rightarrow \mathbb{R}$, where $h(\bm{a}) = 1$ for all column vectors $\bm{a}$ of $A$. A Gale dual of $A$, is an integer $d\times N$ matrix whose rows form a
$\mathbb{Z}$-basis of $L$, we denote this matrix by $B$.

It turns out that $\Phi_{\bm{\gamma}}^L$ satisfies a system of partial differential equations. First of all, let $\bm{m}=
(m_1,\dots,m_N)$ be an integer row vector such that $\bm{m}\cdot\bm{l}=0$ for all $\bm{l}\in L$.
Then, for any $\lambda \in \mathbb{C}^*$, one easily sees that
\begin{gather*}
\Phi_{\bm{\gamma}}^L(\lambda^{m_1}z_1,\dots,\lambda^{m_N}z_N)=
\lambda^{\bm{m}\cdot\bm{\gamma}}\Phi_{\bm{\gamma}}^L(z_1,\dots,z_N).
\end{gather*}
Take the derivative with respect to $\lambda$ and set $\lambda=1$. Then we see that
$\Phi_{\bm{\gamma}}^L$ is annihilated by the differential operator
\begin{gather*}
m_1z_1\partial_{z_1}+\cdots+m_Nz_N\partial{z_N}-\bm{m}\cdot\bm{\gamma}.
\end{gather*}
In particular, if we let $\bm{m}$ be the $i$-th row of $A=(A_{ij})$ we see
that $\Phi_{\bm{\gamma}}^L$ is annihilated by the Euler operator
\begin{gather*}
Z_i:=A_{i1}z_1\partial_{z_1}+\cdots+A_{iN}z_N\partial{z_N}-\alpha_i.
\end{gather*}
There is a second set of differential equations which arises from the
observation
\begin{gather*}
\partial_{z_1}^{\lambda_1}\cdots\partial_{z_N}^{\lambda_N}\Phi_{\bm{\gamma}}^L
=\Phi_{\bm{\gamma}-\bm{\lambda}}^L
\end{gather*}
for any $\bm{\lambda}=(\lambda_1,\dots,\lambda_N)\in\mathbb{Z}_{\ge0}^N$.
Let now $\bm{\lambda}\in L$ and write $\bm{\lambda}=\bm{\lambda}^+-\bm{\lambda}^-$,
where $\bm{\lambda}^\pm$ are integer vectors with non-negative entries. Then,
\begin{gather*}
\partial^{\bm{\lambda}^+}\Phi_{\bm{\gamma}}^L
=\Phi_{\bm{\gamma}-\bm{\lambda}^+}^L=\Phi_{\bm{\gamma}-\bm{\lambda}^-}^L
=\partial^{\bm{\lambda}^-}\Phi_{\bm{\gamma}}^L.
\end{gather*}
We use the notation $\partial^{\bm{\lambda}}=\partial_{z_1}^{\lambda_1}
\cdots\partial_{z_N}^{\lambda_N}$ and the second step follows from the invariance
of $\Phi_{\bm{\gamma}}^L$ when $\bm{\gamma}$ is shifted over vectors in $L$.
Thus we find that $\Phi_{\bm{\gamma}}^L$ is annihilated by the so-called box operators
\begin{gather*}
\Box^{\bm \lambda}:=\prod_{\lambda_i>0}\partial_{z_i}^{\lambda_i}-\prod_{\lambda_i<0}
\partial_{z_i}^{-\lambda_i}
\end{gather*}
for all $\bm{\lambda}\in L$.

{\samepage
The {\it $A$-hypergeometric system} $H_{A}(\bm{\alpha})$ is the system of differential equations generated by
\begin{enumerate}\itemsep=0pt
\item The Euler operators
\begin{gather*}%\label{eq:euler}
Z_j=A_{j1}\partial_{z_1}+\cdots+A_{jN}\partial_{z_N}-\alpha_j,\qquad j=1,\dots,N-d.
\end{gather*}
\item The box operators
\begin{gather*}%\label{eq:box}
\Box^{\bm{\lambda}}=\partial^{\bm{\lambda}^+}-\partial^{\bm{\lambda}^-},\qquad
\bm{\lambda}\in L.
\end{gather*}
\end{enumerate}}

An $A$-hypergeometric function is a holomorphic
function in $z_1,\dots,z_N$ which satisfies the equations in the $A$-hypergeometric system. Either $A$ together with a parameter vector $\bm{\alpha}$ or~$B$
with $\bm{\gamma}$ (modulo translations in $L$) is enough to encode all the information about the $A$-hypergeometric system. The columns of $A$ are denoted by $\bm{a}_1,\dots,\bm{a}_N$ and the columns of $B$ are denoted by $\bm{b}_1,\dots,\bm{b}_N$.
